\section{模型假设}\label{sec:model-assumptions}

在题设设备参数与动作规则的基础上，明确本文模型采用的以下工作条件。其中，全向接收与静默推断的条件仅用于第二、三问；第四问的附加条件见第 7、8 条。

\begin{enumerate}
  \item \textbf{平面运动与坐标准确性。}按题设忽略设备与干扰源的高度差，将两者表示为平面点。假定机器狗能够准确获得自身位置并到达指定坐标，两点间按直线匀速移动，不另计位置误差、地形绕行、加减速与转向代价，使移动耗时可由欧氏距离计算。

  \item \textbf{源参数稳定且各源互不影响。}假定干扰源在一次任务中保持静止，未清除前其频道、发射功率及有效接收半径不变。各源信号按题设互不影响，清除一个源不改变其他源的接收条件。源位置、具体数量及接收半径的真实值均不作为决策输入，只利用公开范围与真实反馈逐步约束。

  \item \textbf{全向信号遵循距离门限。}对于第二、三问中的未清除全向源，是否接收到信号仅由检测距离与该源固定接收半径决定，不额外引入遮挡、间歇发射或随机漏检。距离不超过 $5\,\mathrm{m}$ 时按题设返回近场反馈而非示向度；该条件支持保证接收域、发现覆盖及历史静默推断的建立。

  \item \textbf{测角误差按有界不确定性处理。}保留题设 $\pm1^\circ$ 的误差界，不假定不同地点的误差相互独立或服从特定概率分布，也不通过同地点重复检测缩小误差界。实现中对保留两位小数的示向度采用 $1.005^\circ$ 计算半角，仅用于容纳舍入余量，区域更新仍以真实反馈为依据。

  \item \textbf{动作顺序执行并按规则计费。}移动、换频、检测、光学定位与清除按动作顺序执行，任务总耗时为实际动作费用之和。光学尝试失败时仍计定位耗时，成功后再计清除耗时；算法计算时间另行记录，不与机器狗的动作计费用时混合。

  \item \textbf{光学操作遵循题设距离条件。}假定设备正常工作，在距目标源不超过 $20\,\mathrm{m}$ 时能够完成光学精确定位和清除，其成功条件不依赖信号辐射方向。算法仅在收到真实成功反馈后登记清除；有限步全清的保证以时间和动作预算足够为前提，超时或未确认的请求不作为成功。

  \item \textbf{定向源的辐射范围为半平面。}第四问中的定向源只向发射方向两侧各 $90^\circ$ 的范围辐射，即只有位于过源位置、垂直于发射方向的直线所限定的闭半平面内，且距离不超过接收半径的位置才能收到信号；发射方向在一次任务中固定，两类源的数量和比例均未知。近场反馈同样受该范围限制，光学定位和清除只取决于距离。

  \item \textbf{机器狗可在目标区域外侧附近停留检测。}题目只规定干扰源位于半径 $1800$ m 的目标区域内，未限制机器狗的活动范围；第四问的检测站中有一环位于半径 $1865$ m 的圆周上，用于发现贴着边界朝外发射的定向源。
\end{enumerate}

